\section{Conclusion}

The ansatz
\[
\nabla_{\bm{x}}\times\nabla_{\bm{p}}
=
\frac{\hbar^2}{2\pi}
\]
cannot be repaired by replacing \(\hbar^2/(2\pi)\) with either zero or \(h\). The left-hand side is a generally nonzero, non-intrinsic, vector-valued differential operator, while the proposed right-hand side is a scalar with incompatible dimensions.

The correct hierarchy is
\begin{equation}
\boxed{
\comm{\partial_{x^i}}{\partial_{p_j}}=0,
}
\end{equation}
\begin{equation}
\boxed{
\Pi=\partial_{x^i}\wedge\partial_{p_i},
\qquad
\pb{x^i}{p_j}=\delta^i{}_j,
}
\end{equation}
\begin{equation}
\boxed{
\comm{\hat x^i}{\hat p_j}
=
\starcomm{x^i}{p_j}
=
\ii\hbar\delta^i{}_j,
}
\end{equation}
and, in a separate topological setting,
\begin{equation}
\boxed{
\oint_C\bm{p}\cdot\dd\bm{\ell}=nh.
}
\end{equation}
These equations belong to different levels of mathematical description. Keeping those levels separate preserves dimensional consistency, covariance, and physical meaning while retaining the valid intuition that antisymmetric phase-space structure is central to quantization.

\section*{Editorial classification}

This manuscript is best presented as a pedagogical technical note, mathematical-physics clarification, or foundational appendix. It should not be advertised as a newly discovered quantum identity. For APS-oriented preparation, Physics Subject Headings (PhySH), rather than author-supplied PACS numbers, should be used where required by the venue \cite{APSPhySH}.
